\section{A finite Hankel selection and its interior approximation}
\label{sec:hankel-selection}
This section constructs a complex-place concentration cutoff without
using a zero-location theorem and proves its archimedean approximation
estimates. It does not estimate the inverse Gram matrix of its Hecke
theta images.

Fix $c>0$ and use the unitary involution $\mathcal H_c$ of
Proposition~\ref{prop:complex-concentration}. On $L^2((0,\infty),dx)$
put
\[
 P_\Lambda=1_{(0,\Lambda)},\qquad
 B_\Lambda=\mathcal H_cP_\Lambda\mathcal H_c,\qquad
 A_\Lambda=P_\Lambda B_\Lambda P_\Lambda,
 \qquad \Pi_\Lambda=1_{[1/2,1]}(A_\Lambda).
\]
The last expression is the spectral projection. It defines a finite
Hankel, or circular prolate, selection by concentration, rather than
by an arbitrarily rounded index. It depends on $c$ and $\Lambda$ and
not on the zeros of an $L$-function.
Let $(Uf)(t)=e^{t/2}f(e^t)$, and write $H_D^s$ for the pullback
of the ordinary Sobolev space $H^s(\R,dt)$ by $U$.

\begin{proposition}[Size of the Hankel selection]
\label{prop:hankel-selection-size}
The selected concentration eigenvalues are at least $1/2$, while
every omitted positive eigenvalue is less than $1/2$. With $c$ fixed,
\[
 \operatorname{rank}\Pi_\Lambda
     =\frac{2c\Lambda}{\pi}+O_c(\log(2+\Lambda)).
\]
More precisely,
\[
 \Tr(A_\Lambda-A_\Lambda^2)=O_c(\log(2+\Lambda)),
 \qquad
 |\operatorname{rank}\Pi_\Lambda-\Tr A_\Lambda|
          \leq2\Tr(A_\Lambda-A_\Lambda^2).
\]
\end{proposition}
\begin{proof}
The concentration operator is positive, compact, and at most the
identity. Its trace is finite by
Proposition~\ref{prop:complex-concentration}, so the spectral
selection is finite dimensional. The assertions about selected and
omitted eigenvalues follow from its definition.

Here is a direct estimate of the transition mass. Change variables
$r=2c\sqrt{\Lambda x}$ by the unitary map
\[
 (J_\Lambda f)(r)
    =\left(\frac{r}{2c^2\Lambda}\right)^{1/2}
                        f\left(\frac{r^2}{4c^2\Lambda}\right).
\]
Under this map $B_\Lambda$ becomes the Bessel projection with kernel
\[
 K(r,s)=\sqrt{rs}\int_0^1vJ_0(rv)J_0(sv)\,dv,
\]
and $P_\Lambda$ becomes restriction to $(0,Z)$, where $Z=2c\Lambda$.
The Bessel differential equation, integrated once, gives for $r\ne s$
\[
 K(r,s)=\frac{\sqrt{rs}}{s^2-r^2}
       \bigl(sJ_1(s)J_0(r)-rJ_1(r)J_0(s)\bigr),
 \qquad
 K(r,r)=\frac r2\bigl(J_0(r)^2+J_1(r)^2\bigr).
\]
The standard bounds $|J_0(t)|+|J_1(t)|\leq C(1+t)^{-1/2}$
follow, for example, by splitting their trigonometric integral
representations at distance $t^{-1/2}$ from the stationary points
and integrating by parts on the complement. Their power series
give $J_0(t)=O(1)$ and $J_1(t)=O(t)$ near zero.
For $r,s\geq1$, the displayed off-diagonal formula gives
$|K(r,s)|\leq C|r-s|^{-1}$, while Cauchy--Schwarz in its integral
formula and the diagonal bound give $|K(r,s)|\leq C$.
For $0<r<1$ and $s\geq2$, the same formula gives
$|K(r,s)|\leq C\sqrt r/s$.

Since $B_\Lambda$ is a projection, its kernel therefore yields,
for $Z\geq2$,
\[
 \begin{split}
 \Tr(A_\Lambda-A_\Lambda^2)
 &=\int_0^Z\int_Z^\infty|K(r,s)|^2\,ds\,dr\\
 &\leq C\int_1^Z\int_Z^\infty
          \min(1,|s-r|^{-2})\,ds\,dr+C/Z\\
 &\leq C\log(2+Z).
 \end{split}
\]
The final inequality follows by putting $u=Z-r$ and $v=s-Z$:
the inner integral is at most $C/(1+u)$.
For each $0\leq\nu\leq1$,
$|1_{\nu\geq1/2}-\nu|\leq2\nu(1-\nu)$.
Sum over the eigenvalues and use
$\Tr A_\Lambda=Z/\pi+O(\log(2+Z))$ from
\eqref{eq:hankel-trace}. This proves all assertions.
\end{proof}

\begin{proposition}[Band extensions and dilation regularity]
\label{prop:hankel-raw-estimates}
There is $d>0$ such that, for all sufficiently large $\Lambda$ and
every $\psi\in\operatorname{Ran}\Pi_\Lambda$,
\[
 \int_{d/(c^2\Lambda)}^\Lambda|\psi(x)|^2dx
       \geq\tfrac12\|\psi\|_2^2,
 \qquad
 \|\psi\|_{H_D^s}\leq C_{c,s}(1+\Lambda)^s\|\psi\|_2,
                   \quad0<s<1/2.
\]
The band extension
$W_\Lambda\psi=B_\Lambda(A_\Lambda|_{\operatorname{Ran}\Pi_\Lambda})^{-1}\psi$
satisfies $P_\Lambda W_\Lambda\psi=\psi$ and
$\|W_\Lambda\|\leq\sqrt2$.
Its ordinary raw derivatives obey
\[
 \sup_{x\geq0}|(W_\Lambda\psi)^{(j)}(x)|
             \leq C_{c,j}\Lambda^{j+1/2}\|\psi\|_2
                 \qquad(j=0,1,\ldots).
\]
\end{proposition}
\begin{proof}
On the selected space the inverse of $A_\Lambda$ has norm at most
two, and
\[
 \|W_\Lambda\psi\|_2^2
     =\langle\psi,A_\Lambda^{-1}\psi\rangle\leq2\|\psi\|_2^2.
\]
The equality after restriction is the definition of $A_\Lambda$.
Write $w=W_\Lambda\psi$. Its Hankel transform is supported in
$(0,\Lambda)$. The bound $|J_0|\leq1$ and Cauchy--Schwarz give
$|w(x)|\leq c\sqrt\Lambda\|w\|_2$.
Thus the squared mass in $(0,d/(c^2\Lambda))$ is at most
$2d\|\psi\|_2^2$. Choose $d\leq1/4$.
For the derivatives, the identity
\[
 \partial_x^jJ_0(2c\sqrt{xy})
       =(-c)^j(y/x)^{j/2}J_j(2c\sqrt{xy})
\]
and $|J_j(t)|\leq(t/2)^j/j!$ imply
$|\partial_x^jJ_0(2c\sqrt{xy})|\leq c^{2j}y^j/j!$.
The latter Bessel bound follows directly from its Poisson integral
representation. Integrate the squared bound over $0<y<\Lambda$
and apply Cauchy--Schwarz. The formula extends continuously to
$x=0$ by the Bessel series.

To estimate the dilation derivative, put $F(z)=w(|z|^2)$ on
$\R^2$. The radial Fourier formula shows that $\widehat F$ is
supported in $|\xi|\leq(c/\pi)\sqrt\Lambda$, and
$\|F\|_2^2=\pi\|w\|_2^2$. Plancherel gives
$\|\nabla F\|_2\leq2c\sqrt\Lambda\|F\|_2$.
Since $x\partial_xw$ corresponds to $\tfrac12z\cdot\nabla F$,
\[
 \|xw'(x)\|_{L^2(0,\Lambda)}
             \leq c\Lambda\|w\|_2.
\]
Consequently $f=U\psi$, on $(-\infty,L)$ with $L=\log\Lambda$,
has interior derivative norm at most
$C_c(1+\Lambda)\|\psi\|_2$. Its value tends to zero at the lower
end, since $w$ is bounded near zero. The elementary trace estimate
$|f(L-)|^2\leq2\|f\|_2\|f'\|_2$ and the distributional Fourier
identity for its zero extension show that, for $M\geq1$,
\[
 \|f\|_{H^s}^2\leq C_s\bigl(
 M^{2s}\|f\|_2^2+M^{2s-2}\|f'\|_2^2
                         +M^{2s-1}|f(L-)|^2\bigr).
\]
The last frequency integral converges exactly when $s<1/2$.
Take $M=1+\Lambda$ to obtain the stated estimate.
\end{proof}

\begin{proposition}[Strict-margin Hankel approximation]
\label{prop:hankel-interior-approximation}
Let $p\in C_c^\infty([0,1))$ equal one near zero, with
$0\leq p\leq1$, and set
\[
 P^p_\Lambda=p(x/\Lambda),\qquad
 B^p_\Lambda=\mathcal H_cp(x/\Lambda)\mathcal H_c.
\]
For every $0\leq\sigma<1/2$ and $A>0$,
\[
 \|(1-\Pi_\Lambda)P^p_\Lambda B^p_\Lambda P^p_\Lambda
                         \|_{L^2\to H_D^\sigma}
                   \leq C_{c,p,\sigma,A}\Lambda^{-A}.
\]
The same conclusion holds for each fixed smooth profile supported
strictly inside $(0,1)$, dilated in the norm variable by $\Lambda$,
and for the Hankel transforms of such dilates. For a family already
satisfying $I\psi=\psi(0)=0$, with the stated superpolynomial
projection error and polynomial ordinary local derivative bounds,
the principal constraints can be restored with the same arbitrarily
strong rate, using two projected profiles.
\end{proposition}
\begin{proof}
On radial $L^2(\R^2)$ the two soft operators are multiplication by
$p(|z|^2/\Lambda)$ and the Fourier multiplier
$p((\pi/c)^2|\xi|^2/\Lambda)$. Choose $b<1$ containing the support
of $p$. The latter multiplier is supported in the disk of radius
$(c/\pi)\sqrt{b\Lambda}$, while the hard band has radius
$(c/\pi)\sqrt\Lambda$. The Fourier transform of the former
multiplier is
$\Lambda\widehat{p(|\cdot|^2)}(\sqrt\Lambda\,\xi)$.
Its Schwartz decay and the fixed relative gap between these disks
give, by Schur's test,
\[
 \|(1-B_\Lambda)P^p_\Lambda B^p_\Lambda\|_{2\to2}
                                  =O_{c,p,M}(\Lambda^{-M})
\]
for every fixed $M$. The final soft time multiplication gives
zero hard time tail.
For a vector $f$ supported in $(0,\Lambda)$, the omitted-eigenvalue
bound gives directly
\[
 \|(1-\Pi_\Lambda)f\|_2^2
       \leq2\langle f,(1-A_\Lambda)f\rangle
       =2\|(1-B_\Lambda)f\|_2^2.
\]
This proves the raw superpolynomial estimate.

For any fixed $s<1/2$, the preceding proposition bounds the selected
part in $H_D^s$ by $C\Lambda^s$ times the raw input norm. The soft
part has the same bound: its first dilation derivative costs
$O(1+\Lambda)$ by the radial band estimate, and it vanishes before
the hard endpoint. Interpolate this coarse $H_D^s$ bound with the
raw estimate, taking $\sigma<s<1/2$ and then taking $M$ large.
This proves the operator estimate for every $A$.

Let $(D_r f)(x)=r^{-1/2}f(x/r)$. Direct substitution gives
$\mathcal H_cD_r=D_{1/r}\mathcal H_c$, so the Hankel transform
reverses the dilation generator and preserves every $H_D^s$ norm.
For a fixed smooth profile $\beta$ in the statement,
$\mathcal H_c\beta$ is Schwartz in the radial variable, hence has
every polynomial decay rate in the norm variable. The identity
$\mathcal H_cD_\Lambda\beta=D_{1/\Lambda}\mathcal H_c\beta$
therefore gives every polynomially small hard band tail, while
the raw time tail is zero. For its Hankel dual the roles reverse.
For a profile with a small raw time tail, first restrict it by
$P_\Lambda$; the resulting band tail increases by at most that time
tail. Thus the preceding projection bound is at most a fixed
constant times the sum of both raw tails. The same interpolation
proves the claimed fixed-profile conclusions.

Choose $\beta$ supported away from zero with $I\beta=1$.
The two profiles
\[
 b_{\Lambda,+}=D_\Lambda\beta,\qquad
 b_{\Lambda,-}=c^{-1}\mathcal H_cD_\Lambda\beta
\]
have identity matrix for the normalized constraints
$I/\sqrt\Lambda$ and $\operatorname{ev}_0/\sqrt\Lambda$:
indeed $(\mathcal H_cf)(0)=cIf$ and
$I\mathcal H_cf=f(0)/c$ for these Schwartz profiles.
Their projected constraint matrix tends to the identity with
every polynomial rate. For the mean, use Cauchy--Schwarz on the
time interval and the Schwartz exterior tail of the dual profile.
For evaluation, the projection error has every raw $L^2$ rate,
whereas its ordinary derivative on $(0,1)$ has polynomial growth
by the preceding band-extension estimate. The one-dimensional
trace inequality
\[
 |F(0)|^2\leq C\|F\|_{L^2(0,1)}
          \bigl(\|F'\|_{L^2(0,1)}+\|F\|_{L^2(0,1)}\bigr)
\]
proves the evaluation assertion. Subtracting the two projected
profiles multiplied by the inverse constraint matrix restores
the constraints exactly. The correction has every prescribed
$H_D^\sigma$ rate whenever the original constraint errors do.
\end{proof}

If a Hecke theta multiplier satisfies a fixed bound
$|L_K(1/2+i\tau,\chi)|\leq C(1+|\tau|)^\sigma$ for some
$\sigma<1/2$, the last proposition immediately gives theta graph
approximation, by Plancherel. Such a subconvex estimate is a
separate arithmetic input. Neither that estimate nor the purely
archimedean propositions prove the inverse theta Gram inequality
or the smoothed theta-tail estimate needed for a Hecke trace limit.
Proposition~\ref{prop:hecke-crude-gram},
Theorem~\ref{thm:hecke-deep-leakage}, and
Theorem~\ref{thm:hi-local-limit} establish those arithmetic and
local-limit comparisons for the deeper selection below.
